\chapter{Discusión}\label{cap:discusion}

\section{Qué aporta modelar la demanda como un MBG}

La diferencia entre evaluar un diseño en la demanda esperada y evaluarlo sobre miles de escenarios no es
cosmética. Con congestión de tipo BPR, el costo social es convexo en la demanda, y la asimetría de las
bandas de la \cref{fig:trayectorias} lo muestra: los escenarios de alta demanda empeoran el sistema
mucho más de lo que lo alivian los de baja demanda. La fase I añade una lección complementaria: la
demanda que efectivamente ocurrió entre 2013 y 2025 quedó en el percentil \PercentilObsFin{} de los
escenarios construidos con la información de 2013. La cola existía en la distribución, pero un
análisis con un escenario central y una sensibilidad de más o menos 10\,\% no la habría contemplado.
La extensión con saltos de tipo Merton existe precisamente para dar más peso a esos eventos, y con
ella las líneas de metro, que reducen la dispersión del costo, ganan atractivo frente a criterios
centrados en la media.

\section{La tasa de descuento como variable de decisión}

En ambas fases la tasa social de descuento es el supuesto que más cambia el orden de los diseños. En la
fase II, a 9\,\% las líneas de metro adicionales no se pagan, y a 3,5\,\% todas lo hacen. En la fase I,
la tasa decide entre elevado y subterráneo: a 12\,\% gana el elevado, a 9\,\% empatan con el costo
oficial y por debajo de \TasaEmpate{} gana el subterráneo. Esto no es un defecto del modelo sino la
expresión de un juicio de valor: cuánto pesa el bienestar de quienes usarán el metro en 2050 frente al
costo que se paga hoy. La recomendación práctica es que las decisiones de tipología se presenten
siempre con la tasa de quiebre explícita, para que el debate se haga sobre ese juicio y no sobre cifras
aisladas.

\section{El debate sobre el metro elevado}

El debate público sobre la Primera Línea se ha planteado a menudo como una cuestión de principio
(«una ciudad de este tamaño merece un metro subterráneo»). El modelo sugiere reformularlo como tres
preguntas cuantificables: cuál es el costo creíble del subterráneo, qué tasa de descuento se usa y
cuánto vale el costo urbano del viaducto. Con las respuestas de la época (costo oficial, 12\,\% y
pérdida de valor de los predios como medida del costo urbano), el elevado era preferible. Con otras
respuestas igualmente defendibles (9\,\% y un costo urbano en el extremo alto del rango), el
subterráneo de 2014 era competitivo. Lo que el modelo no respalda es que el tamaño de la ciudad baste:
el tamaño exige una red, y la red favorece los kilómetros más baratos.

La fase I muestra, además, que el mayor costo de la Primera Línea no vino de la tipología sino del
tiempo: entre el diseño de 2014 y la operación prevista en 2028 pasaron catorce años, y en ese lapso la
demanda resultó menor de lo esperado. Si algo enseña la historia del metro de Bogotá es que cambiar de
decisión cuesta, y que el valor de una decisión de infraestructura depende tanto de la tipología como
de la capacidad de ejecutarla a tiempo.

\section{Limitaciones}

\begin{enumerate}
\item \textbf{Escala.} El modelo tiene doce zonas y dieciocho corredores. Sirve para comparar diseños
      estratégicos con los mismos escenarios, no para dimensionar estaciones ni prever cargas por tramo.
\item \textbf{Ubicación y tipología.} Las líneas se representan como corredores entre zonas. El modelo
      no distingue entre tramos con avenidas anchas, donde un viaducto es viable, y tejidos estrechos,
      donde exigiría compras de predios; las variantes elevadas de las líneas 2 y 3 suponen que el
      viaducto es viable.
\item \textbf{Costo urbano del viaducto.} Se trata mediante umbrales y una estimación externa basada en
      el valor catastral, no dentro del modelo. Una estimación hedónica con datos de Bogotá, por ejemplo
      con las transacciones alrededor del viaducto cuando entre en operación, permitiría cerrar esta
      brecha.
\item \textbf{Sobrecostos.} Las dos tipologías comparten la misma distribución de sobrecostos. Si los
      túneles en los suelos blandos de la Sabana tuvieran más riesgo, el subterráneo saldría peor.
\item \textbf{Trayectoria observada.} La trayectoria 2013--2025 de la fase I es estilizada: combina
      la evolución de la demanda de TransMilenio con las encuestas de movilidad.
\item \textbf{Uso del suelo.} La política de proximidad se representa con dos efectos agregados
      (menos viajes motorizados y relocalización de la atracción) y un costo anual supuesto. Sus
      beneficios dependen de que la política efectivamente cambie la localización de empleos y
      viviendas, algo que el modelo supone y no demuestra.
\end{enumerate}

\section{Extensiones}

Tres extensiones parecen las más valiosas. La primera es endogenizar el uso del suelo con un modelo de
localización que responda a la accesibilidad, en la línea de \citet{alonso1964} y \citet{muth1969}.
La segunda es estimar el costo urbano de los viaductos con datos de Bogotá cuando la Primera Línea
entre en operación. La tercera es valorar como opciones reales no solo la Línea 3, sino la secuencia
completa de la red, incluida la opción de soterrar tramos futuros si la demanda o el costo urbano lo
justifican.
